% !TEX root = ../main.tex
% =====================================================================
% 7. FINDING 3 -- AOPCR is not comparable across configurations.
% Budget 0.75 pp.
%
% This is the most original section in the paper. It is a criticism of
% an evaluation practice, backed by a controlled measurement, and it
% must be written carefully: the point is about the METRIC's scale, not
% about anyone's competence.
% =====================================================================

\section{AOPCR Is Not Comparable Across Training Configurations}
\label{sec:aopcr}

\begin{finding}
\label{find:aopcr}
Holding the architecture, the code, the metric implementation and the data fixed, and changing only
the training recipe, moves AOPCR by a factor of five. AOPCR therefore cannot support comparisons
between models trained under different configurations -- which is what a comparison across papers
is.
\end{finding}

\subsection{The measurement}
\label{sec:aopcr:measure}

The comparison in \Cref{tab:aopcr} is unusually clean, because both rows are literally the same
model. Same InceptionTime encoder, same conjunctive head, same parameter count, same data loaders,
same AOPCR implementation, evaluated in the same harness on the same machine. The only difference is
the training recipe: our $400$-epoch configuration of \Cref{sec:protocol} against \millet{}'s own
published $1500$-epoch recipe.

\begin{table}[t]
\caption{The same architecture and the same code under two training recipes. Accuracy moves by a few
points, as expected from a longer schedule. AOPCR moves by a factor of five, on both datasets.}
\label{tab:aopcr}
\begin{center}
\small
\begin{tabular}{lcccc}
\toprule
& \multicolumn{2}{c}{\textbf{\webtraffic{}}} & \multicolumn{2}{c}{\textbf{\ucr{}-85}} \\
\cmidrule(lr){2-3}\cmidrule(lr){4-5}
\textbf{Training recipe} & \textbf{Accuracy} & \textbf{AOPCR} & \textbf{Accuracy} & \textbf{AOPCR} \\
\midrule
Ours, $400$ epochs (\Cref{sec:protocol})   & 0.887 & \phantom{0}2.569 & 0.8274 & 0.8119 \\
\millet{}'s, $1500$ epochs                  & 0.920 & \textbf{13.268} & 0.8434 & \textbf{3.9264} \\
\midrule
\emph{ratio}                                & $1.04\times$ & $\mathbf{5.17\times}$ & $1.02\times$ & $\mathbf{4.84\times}$ \\
\bottomrule
\end{tabular}
\end{center}
\end{table}

A $5\times$ change in a metric that is supposed to measure explanation quality, produced by a change
that improved accuracy by three points, is not a $5\times$ improvement in explanation quality.

\subsection{Why this happens}
\label{sec:aopcr:why}

The cause is visible in \Cref{eq:aopc}. AOPC accumulates \emph{raw confidence differences}
$\bagpred^{k}(\bag) - \bagpred^{k}(\bag^{O^k}_{j})$, and AOPCR takes the margin of that quantity over
random maskings. Neither operation normalises by the range the confidence can move over. So the
metric inherits the scale of the model's outputs, and anything that changes that scale changes AOPCR
without changing the ranking it is meant to evaluate. A longer schedule with no weight decay and no
label smoothing -- \millet{}'s recipe -- produces exactly that: larger, more saturated logits, hence
larger differences when a step is masked, hence a larger area. The ordering of the time steps, which
is the thing an explanation actually asserts, need not have improved at all.

This also explains the behaviour noted in \Cref{sec:selection:kappa}, where AOPCR failed to track
$\kappa$ while \ndcg{} did. A metric whose scale is set by the confidence distribution will be a
noisy instrument for changes that act on the ranking.

\subsection{What follows for practice}
\label{sec:aopcr:protocol}

We are not arguing that AOPCR should be abandoned. It is the only faithfulness measure available on
datasets without per-step ground truth, which is nearly all of them, and \emph{within} a fixed
configuration it does its job. We argue that three things should accompany it, none of which costs
anything:

\begin{enumerate}
  \item \textbf{Report the training configuration next to the AOPCR}, and compare AOPCR only within
        one. A table of AOPCR values whose rows come from different papers has no defined meaning.
  \item \textbf{Re-train the baseline rather than quoting it}, when the claim involves a
        faithfulness metric. This is the reason every baseline in this paper was re-trained, and
        \Cref{tab:aopcr} is the evidence that the extra compute was not optional.
  \item \textbf{Prefer a normalised metric when ground truth exists.} \ndcg{} is bounded in $[0,1]$
        and is unaffected by logit scale; where the two disagree it is the stronger evidence.
\end{enumerate}

Applied to ourselves, this is why \Cref{tab:web} carries a warning on its AOPCR column, and why we
never place any AOPCR in this paper beside the value \millet{} published. Applied to the literature,
it means that published AOPCR comparisons between independently trained models -- including
comparisons we could have made in our own favour -- should be read as uninformative rather than as
evidence. A normalised replacement that preserves AOPCR's ground-truth-free property is the natural
next step, and we did not have the compute budget to validate one here.
